% ============================================================================
%  Merged implementation chapter:
%  implementation_multi_sc.tex + implementation_two_finger.tex
% ============================================================================

This chapter presents the optimization frameworks developed for two gripper types: a
multi-suction cup gripper consisting of three suction cups arranged in a triangular configuration,
and a two-finger gripper with a spline-based gripping contour. Both frameworks support evolutionary
multi-objective optimization and Bayesian multi-objective optimization. The evolutionary algorithms
are implemented using the \texttt{pymoo} library~\cite{pymoo}, whereas Bayesian optimization is
implemented using the \texttt{Ax} platform and its underlying \texttt{BoTorch} modeling
framework~\cite{olson2025ax,botorch}.

For both gripper types and both optimization approaches, the evaluation of candidate gripper
designs is delegated to an external service through an API. This separation provides a modular
architecture in which the optimization strategy and the domain-specific grasp evaluation remain
decoupled. Consequently, the same decision-variable representation, geometry-generation pipeline,
grasp-generation procedure, and objective functions can be used by both evolutionary and Bayesian
optimization methods, allowing their performance to be compared under a common evaluation budget
and a common evaluation pipeline.

\section{Implementation of the Common Gripper Optimization Framework}
\label{sec:gripper_optimization_framework}

This section describes the parts of the implementation that are shared by both problems: the
configuration, the supported optimization methods, the initialization, the optimization loops, the
termination criteria, the performance monitoring, the optional hyperparameter optimization, the
objective-direction conventions, and the output. The problem-specific parameterizations, geometry
pipelines, and objective definitions are described in
Section~\ref{sec:multi_sc_gripper_optimization} for the multi-suction cup gripper and in
Section~\ref{sec:two_finger_gripper_optimization} for the two-finger gripper. An overview of the
two problems is given in Table~\ref{tab:problem_overview}.

\begin{table}[htbp]
\centering
\caption{Overview of the two gripper optimization problems.}
\label{tab:problem_overview}
\begin{tabularx}{\textwidth}{@{}l >{\raggedright\arraybackslash}X >{\raggedright\arraybackslash}X@{}}
\toprule
 & \textbf{Multi-suction cup gripper} & \textbf{Two-finger gripper} \\
\midrule
Decision vector &
$\mathbf{q}=[r,\alpha_1,\alpha_2,\alpha_3]$ &
$\mathbf{q}=[\ell_{\mathrm{gn}},h_1,h_2,h_3]^{\mathrm{T}}$ \\

Dimension $d$ & 4 & 4 \\

Objectives &
$m=3$: diversity $D$, normalised grasp count $\hat{n}_{\mathrm{g}}$,
normalised triangle area $\hat{A}$ &
$m=n_{\mathrm{or}}=10$: one proximity KPI $k_i$ per predefined workpiece orientation \\

Evaluation core &
Suction grasp generation, seal-formation and collision filtering, clustering &
CadQuery geometry generation, voxelization, FFT-based collision and proximity analysis \\
\bottomrule
\end{tabularx}
\end{table}

\subsection{Optimization Framework Configuration}
\label{sec:optimization_framework_sc}
\label{sec:two_finger_framework}

Each gripper problem is solved using either an evolutionary multi-objective optimization algorithm
or Bayesian multi-objective optimization. Both approaches operate on the same decision vector,
variable bounds, objective functions, and external evaluation service.

The common optimization configuration contains:
\begin{itemize}
    \item Number of objectives $m$
    \item Decision-variable definitions and bounds
    \item Optimization direction (minimization or maximization)
    \item Choice of optimization method
    \item Evaluation endpoint URL
\end{itemize}

For evolutionary optimization, the additional configuration parameters include:
\begin{itemize}
    \item Population size $N_{\mathrm{pop}}$
    \item Maximum number of generations $G_{\max}$
    \item Choice of evolutionary algorithm
    \item Algorithm-specific parameters, such as crossover and mutation settings
\end{itemize}

For Bayesian optimization, the additional configuration parameters include:
\begin{itemize}
    \item Maximum number of candidate-generation batches $B_{\max}$, which is read from the same
          configuration field as the number of generations of the evolutionary algorithms
    \item Maximum number of candidates requested per batch $b$
\end{itemize}

In the current Bayesian implementation, the configuration value named \texttt{generations}
represents the number of candidate-generation batches rather than evolutionary generations. A batch
therefore represents a candidate-generation step rather than an evolutionary generation. The maximum
requested number of Bayesian evaluations is
\begin{equation}
    E_{\max}^{\mathrm{BO}} = B_{\max}\,b .
\end{equation}
The actual number of evaluations may be smaller if the generation strategy returns fewer than $b$
candidates in a particular batch.
\nomenclature{$B_{\max}$}{Maximum number of Bayesian optimization batches}
\nomenclature{$b$}{Maximum number of Bayesian candidates requested per batch}
\nomenclature{$E_{\max}^{\mathrm{BO}}$}{Maximum requested number of Bayesian objective evaluations}

\subsubsection{Supported Optimization Methods}

The framework supports the evolutionary and Bayesian multi-objective optimization methods listed in
Table~\ref{tab:algorithms_sc}.

\begin{table}[H]
\centering
\caption{Supported multi-objective optimization methods.}
\label{tab:algorithms_sc}
\label{tab:two_finger_algorithms}
\begin{tabularx}{\textwidth}{@{}l l >{\raggedright\arraybackslash}X@{}}
\toprule
\textbf{Family} & \textbf{Identifier} & \textbf{Method} \\
\midrule
Evolutionary & NSGA2   & Non-dominated Sorting Genetic Algorithm II \\
Evolutionary & NSGA3   & Non-dominated Sorting Genetic Algorithm III \\
Evolutionary & MOEAD   & Multi-Objective Evolutionary Algorithm based on Decomposition \\
Evolutionary & CMOPSO  & Competitive Multi-Objective Particle Swarm Optimization \\
Evolutionary & SMSEMOA & S-Metric Selection Evolutionary Multi-Objective Algorithm \\
Bayesian     & AX-MBM  & Ax modular BoTorch model for multi-objective Bayesian optimization \\
\bottomrule
\end{tabularx}
\end{table}

\subsection{Initialization}
\label{sec:two_finger_initialization}

The initialization procedure depends on the selected optimization approach.

\subsubsection{Evolutionary Initialization}

For evolutionary optimization, an initial population of $N_{\mathrm{pop}}$ candidate solutions is
generated by sampling the decision variables within their bounds:
\begin{equation}
    \begin{gathered}
        \mathcal{U}^{(0)}
        =
        \left\{
            \mathbf{q}^{(0)}_1,
            \mathbf{q}^{(0)}_2,
            \ldots,
            \mathbf{q}^{(0)}_{N_{\mathrm{pop}}}
        \right\}, \\[4pt]
        \mathcal{Q} \subseteq \mathbb{R}^{d},\quad
        d \in \mathbb{N},\quad
        \mathbf{q}^{(0)}_k \in \mathcal{Q},\quad
        k \in \{1,\ldots,N_{\mathrm{pop}}\},\quad
        N_{\mathrm{pop}} \in \mathbb{N},\\[2pt]
        \mathcal{U}^{(0)} \subset \mathcal{Q},\quad
        \bigl|\mathcal{U}^{(0)}\bigr| = N_{\mathrm{pop}}
    \end{gathered}
\end{equation}
\nomenclature{$\mathcal{U}^{(G)}$}{Evolutionary population of candidate decision vectors at generation $G$}
\nomenclature{$N_{\mathrm{pop}}$}{Population size of the evolutionary algorithms}

The complete population is evaluated through the external API before the selected evolutionary
algorithm performs its first selection and variation step. For the multi-suction cup gripper, each
candidate decision vector $\mathbf{q}_k$ is transformed into its corresponding Cartesian gripper
configuration $\mathbf{P}_k$ before objective evaluation, as described in
Section~\ref{sec:multi_sc_api_evaluation}.

\subsubsection{Bayesian Initialization}

Bayesian optimization does not maintain an evolutionary population. Instead, it stores a dataset of
completed trials:
\begin{equation}
    \begin{gathered}
        \mathcal{D}_t =
        \left\{
            \left(\mathbf{q}_k,\mathbf{f}(\mathbf{q}_k)\right)
        \right\},
        \qquad k = 1,2,\ldots,t, \\[4pt]
        \mathbf{q}_k \in \mathcal{Q} \subset \mathbb{R}^{d},\quad
        \mathbf{f} : \mathcal{Q} \rightarrow \mathcal{Y} \subseteq \mathbb{R}^{m},\quad
        m \in \mathbb{N},\; m \geq 2,\quad
        t \in \mathbb{N}_0,\\[2pt]
        \mathcal{D}_t \subset \mathcal{Q} \times \mathcal{Y}
        \;\text{(multiset of } t \text{ trials)},\quad
        \mathcal{D}_0 = \emptyset,\quad
        \mathcal{D}_{t} \subseteq \mathcal{D}_{t+1}
    \end{gathered}
\end{equation}
where $\mathbf{q}_k$ is an evaluated gripper parameterization and $\mathbf{f}(\mathbf{q}_k)$ is its
observed objective vector as returned by the evaluation service.
\nomenclature{$\mathcal{D}_t$}{Dataset of completed Bayesian optimization trials after $t$ evaluations}
\nomenclature{$t$}{Number of completed Bayesian trials}

The Ax generation strategy first evaluates the center of the bounded search space and then generates
an initial set of candidates using Sobol quasi-random sampling. Sobol sampling provides broad and
approximately space-filling coverage of the four-dimensional decision space before a surrogate model
is fitted. Once the initialization requirements of the generation strategy have been satisfied, Ax
transitions to the model-based modular BoTorch model (MBM) stage. From this point onward, new
gripper designs are selected using surrogate predictions and predictive uncertainty.

\subsection{Objective-Direction Conventions}
\label{sec:objective_direction_conventions}

For both gripper problems, all objectives are maximised in their natural form, so that the natural
optimization problem is
\begin{equation}
    \max_{\mathbf{q}\in\mathcal{Q}} \mathbf{f}(\mathbf{q}),
\end{equation}
where $\mathcal{Q}$ is the bounded decision space introduced in
Section~\ref{subsec:methodology_moo}.
\nomenclature{$\mathbf{f}(\cdot)$}{Natural objective vector returned by the evaluation service}

The objective-direction convention is handled differently by the two optimization frameworks. Since
\texttt{pymoo} assumes minimisation, the evolutionary algorithms receive the negated objective
vector
\begin{equation}
    \mathbf{f}_{\mathrm{EA}}(\mathbf{q}) = -\,\mathbf{f}(\mathbf{q}).
\end{equation}
\nomenclature{$\mathbf{f}_{\mathrm{EA}}(\cdot)$}{Minimisation-form objective vector supplied to the evolutionary algorithms}

In the Bayesian implementation, the values returned by the external evaluation service are left
unchanged, since the maximisation direction is declared through the \texttt{Ax} optimization
configuration:
\begin{equation}
    \mathbf{f}_{\mathrm{BO}}(\mathbf{q}) = \mathbf{f}(\mathbf{q}).
\end{equation}
\nomenclature{$\mathbf{f}_{\mathrm{BO}}(\cdot)$}{Objective observations supplied to Ax}

Consequently, the observations are not negated before they are attached to the corresponding Ax
trial, which preserves the original physical meaning of the objective values in the returned Pareto
set. For performance comparisons, the results of both optimization frameworks are converted to the
common natural maximisation convention $\mathbf{f}(\mathbf{q})$, so that all reported objective
values correspond to the original direction and magnitude of the underlying quality measures.

\subsection{Optimization Loop}
\label{sec:two_finger_optimization_loop}

The external evaluation pipeline is shared by the evolutionary and Bayesian optimization
approaches. The methods differ primarily in how they use the completed evaluations to generate
subsequent candidates, while the external candidate-evaluation procedure remains identical.

For every candidate, the common evaluation steps are:
\begin{enumerate}
    \item Send the candidate decision vector to the external evaluation API.
    \item Map the decision variables to the complete problem-specific gripper configuration.
    \item Evaluate the candidate using the problem-specific grasp-evaluation pipeline.
    \item Return the objective vector to the selected optimization framework.
\end{enumerate}
The problem-specific instantiations of these steps are given in
Section~\ref{sec:multi_sc_api_evaluation} and Section~\ref{sec:two_finger_api_evaluation}.

\subsubsection{Evolutionary Optimization Loop}

In every evolutionary generation, the selected algorithm generates a population of candidate
decision vectors. The candidates are evaluated through the common evaluation pipeline, and the
resulting objective values are negated before being supplied to \texttt{pymoo}. The population
update can be expressed abstractly as
\begin{equation}
    \begin{gathered}
        \mathcal{U}^{(G+1)}
        =
        \mathcal{A}
        \left(
            \mathcal{U}^{(G)},
            \mathbf{F}_{\mathrm{EA}}^{(G)}
        \right),
        \qquad
        G = 1,2,\ldots, G_{\max}, \\[4pt]
        \mathcal{U}^{(G)} \subset \mathcal{Q} \subset \mathbb{R}^{d},\quad
        \bigl|\mathcal{U}^{(G)}\bigr| \leq N_{\mathrm{pop}},\quad
        \mathbf{F}_{\mathrm{EA}}^{(G)} \in \mathbb{R}^{n^{(G)} \times m},\quad
        n^{(G)} \in \mathbb{N},\\[2pt]
        \mathcal{A} :
        \mathcal{Q}^{N_{\mathrm{pop}}} \times \mathbb{R}^{n^{(G)} \times m} \times \Omega
        \rightarrow
        \mathcal{Q}^{N_{\mathrm{pop}}},\quad
        G \in \mathbb{N},\quad
        G_{\max} \in \mathbb{N}
    \end{gathered}
\end{equation}
where $\mathcal{A}(\cdot,\cdot)$ denotes the selected evolutionary algorithm, which maps the current
population and the minimisation-form objective values onto the next population, and
$\mathbf{F}_{\mathrm{EA}}^{(G)}$ contains the minimisation-form objective vectors evaluated in
generation $G$.
\nomenclature{$\mathcal{A}(\cdot,\cdot)$}{Selected multi-objective evolutionary algorithm}
\nomenclature{$G$}{Evolutionary generation index}
\nomenclature{$\mathbf{F}_{\mathrm{EA}}^{(G)}$}{Minimisation-form objective values of the population in generation $G$}

Depending on the selected algorithm, the new population is generated through operations such as
selection, crossover, mutation, decomposition, particle movement, or survival based on hypervolume
contribution as explained in Section~\ref{subsubsec:evolutionary_algorithms}. Every candidate in the resulting population is evaluated through the common API-based
evaluation pipeline.

\subsubsection{Bayesian Optimization Loop}

The Bayesian optimization procedure maintains the observation dataset $\mathcal{D}_t$. During the
initialization phase, candidates are generated from the center of the search space and by Sobol
sampling. During the model-based stage, Ax and BoTorch use the completed trial dataset to fit
probabilistic surrogate models:
\begin{equation}
    \begin{gathered}
        \widehat{\mathbf{f}}_t
        =
        \mathcal{M}(\mathcal{D}_t), \\[4pt]
        \mathcal{M} :
        \bigcup_{\tau \geq t_{\mathrm{init}}}
        \bigl(\mathcal{Q}\times\mathcal{Y}\bigr)^{\tau}
        \rightarrow
        \bigl\{\,\widehat{\mathbf{f}} : \mathcal{Q} \rightarrow \Delta(\mathbb{R}^{m})\,\bigr\},
        \qquad t \geq t_{\mathrm{init}},\\[4pt]
        \widehat{\mathbf{f}}_t(\mathbf{q})
        \sim
        \mathcal{N}\!\left(
            \boldsymbol{\mu}_t(\mathbf{q}),\,
            \boldsymbol{\Sigma}_t(\mathbf{q})
        \right),\quad
        \boldsymbol{\mu}_t : \mathcal{Q} \rightarrow \mathbb{R}^{m},\quad
        \boldsymbol{\Sigma}_t : \mathcal{Q} \rightarrow \mathbb{S}^{m}_{+}
    \end{gathered}
\end{equation}
where $\mathcal{M}(\cdot)$ denotes the surrogate-modeling procedure used by Ax and BoTorch.
\nomenclature{$\mathcal{M}(\cdot)$}{Bayesian surrogate-modeling procedure}
\nomenclature{$\widehat{\mathbf{f}}_t$}{Probabilistic surrogate approximation of the objective vector after $t$ evaluations}

The surrogate models approximate the relationship between the gripper parameters and the objective
values. In addition to predicted objective values, they estimate predictive uncertainty in regions
containing few or no observations. An acquisition function $\psi(\cdot\mid\mathcal{D}_t)$ uses this
information to balance:
\begin{itemize}
    \item Exploitation of gripper designs expected to improve the current Pareto frontier
    \item Exploration of uncertain regions of the bounded decision space
\end{itemize}

For a requested batch size $b$, the next candidate batch is generated as
\begin{equation}
    \begin{gathered}
        \mathcal{B}_{t+1}
        =
        \operatorname*{arg\,max}_{
            \substack{
                \mathcal{B}\subset\mathcal{Q}\\
                |\mathcal{B}|\leq b
            }
        }
        \psi(\mathcal{B}\mid\mathcal{D}_t), \\[4pt]
        \psi(\,\cdot \mid \mathcal{D}_t) :
        \bigl\{\mathcal{B} \subset \mathcal{Q} : |\mathcal{B}| \leq b \bigr\}
        \rightarrow \mathbb{R},\quad
        b \in \mathbb{N},\\[2pt]
        \mathcal{B}_{t+1} \subset \mathcal{Q} \subset \mathbb{R}^{d},\quad
        0 \leq |\mathcal{B}_{t+1}| \leq b
        \;\text{(the strategy may return fewer than } b \text{ candidates)}
    \end{gathered}
\end{equation}
\nomenclature{$\mathcal{B}_{t+1}$}{Batch of candidates generated by Bayesian optimization}
\nomenclature{$\psi(\cdot\mid\mathcal{D}_t)$}{Bayesian acquisition function conditioned on the observed trials}

Each candidate $\mathbf{q}\in\mathcal{B}_{t+1}$ is evaluated through the external evaluation
service, and the resulting observations are supplied to Ax as completed trials. After all candidates
of the generated batch have been evaluated, the dataset is updated as
\begin{equation}
    \mathcal{D}_{t+|\mathcal{B}_{t+1}|}
    =
    \mathcal{D}_t
    \cup
    \left\{
        \left(
            \mathbf{q},\,
            \mathbf{f}(\mathbf{q})
        \right)
        \;\middle|\;
        \mathbf{q}\in\mathcal{B}_{t+1}
    \right\}.
\end{equation}

The surrogate models are subsequently updated using the expanded dataset, and another candidate
batch is generated. This process is repeated until the specified evaluation budget is exhausted.

Although multiple candidates can be requested in one Bayesian batch, the current implementation
submits the candidates of that batch one after another to the external evaluation service. The
implementation therefore performs \emph{sequential batch evaluation} rather than parallel batch
evaluation.

\subsection{Termination Criteria}
\label{sec:two_finger_termination}

The termination criteria depend on the selected optimization approach and are summarised in
Table~\ref{tab:termination_sc}.

\begin{table}[htbp]
\centering
\caption{Termination criteria for evolutionary and Bayesian optimization.}
\label{tab:termination_sc}
\label{tab:two_finger_termination}
\begin{tabularx}{\textwidth}{@{}l l l >{\raggedright\arraybackslash}X@{}}
\toprule
\textbf{Approach} & \textbf{Criterion} & \textbf{Symbol} & \textbf{Description} \\
\midrule
Evolutionary & Design-space tolerance & $q_{\text{tol}}$ & Change in decision variables below
threshold \\

Evolutionary & Objective-space tolerance & $f_{\text{tol}}$ & Relative objective improvement below
threshold \\

Evolutionary & Maximum generations & $G_{\max}$ & Maximum number of generations reached \\

Evolutionary & Maximum evaluations & $E_{\max}$ & Maximum number of function evaluations reached \\

Evolutionary & Hypervolume stagnation & $\Delta\mathrm{HV}_{\text{tol}}$ & Hypervolume improvement
over the monitoring window below threshold \\

Bayesian & Maximum batches & $B_{\max}$ & Maximum number of candidate-generation batches reached \\

Bayesian & Maximum evaluations & $E_{\max}^{\mathrm{BO}}$ & Bayesian evaluation budget exhausted \\
\bottomrule
\end{tabularx}
\end{table}

For evolutionary optimization, the design-space and objective-space tolerances are evaluated over a
sliding window of $n_{\mathrm{win}}$ generations.
\nomenclature{$q_{\text{tol}}$}{Design-space convergence tolerance}
\nomenclature{$f_{\text{tol}}$}{Objective-space convergence tolerance}
\nomenclature{$\Delta\mathrm{HV}_{\text{tol}}$}{Hypervolume-stagnation tolerance}
\nomenclature{$n_{\mathrm{win}}$}{Sliding-window length of the evolutionary termination criteria}
\nomenclature{$G_{\max}$}{Maximum number of evolutionary generations}
\nomenclature{$E_{\max}$}{Maximum number of evolutionary objective evaluations}

In the current Bayesian implementation, termination is budget-based. The optimization performs at
most $B_{\max}$ candidate-generation batches, with up to $b$ candidates requested per batch.
Therefore,
\begin{equation}
    E_{\max}^{\mathrm{BO}} = B_{\max}\,b .
\end{equation}
No design-space or objective-space convergence tolerance is currently applied to the Bayesian
optimization loop.

\subsection{Performance Monitoring}
\label{sec:two_finger_monitoring}

The hypervolume indicator is used to monitor both evolutionary and Bayesian optimization. Before
calculating the hypervolume, all objective vectors are represented using the common natural
maximisation convention $\mathbf{f}(\mathbf{q})$ introduced in
Section~\ref{sec:objective_direction_conventions}.

For evolutionary optimization, a callback records the optimization state after each generation, so
that the hypervolume indicator is tracked per generation:
\begin{equation}
    \begin{gathered}
        \mathrm{HV}^{(G)}
        =
        \mathrm{HV}\!\left(\mathcal{F}^{(G)},\mathbf{z}^{\mathrm{ref}}\right), \\[4pt]
        \mathcal{F}^{(G)}
        =
        \left\{\mathbf{f}(\mathbf{q}) \;\middle|\; \mathbf{q}\in\mathcal{U}^{(G)}\right\}
        \subset \mathcal{Y} \subseteq \mathbb{R}^{m},\quad
        \bigl|\mathcal{F}^{(G)}\bigr| \leq N_{\mathrm{pop}},\\[2pt]
        \mathbf{z}^{\mathrm{ref}} \in \mathbb{R}^{m}
        \;\text{with}\;
        z^{\mathrm{ref}}_i \leq \inf_{\mathbf{q}\in\mathcal{Q}} f_i(\mathbf{q}),\;
        i = 1,\ldots,m,\\[2pt]
        \mathrm{HV} :
        \bigl\{\mathcal{F}\subset\mathbb{R}^{m} : |\mathcal{F}|<\infty\bigr\}
        \times \mathbb{R}^{m}
        \rightarrow \mathbb{R}_{\geq 0},\quad
        \mathrm{HV}(\emptyset,\mathbf{z}^{\mathrm{ref}}) = 0
    \end{gathered}
\end{equation}
where $\mathbf{z}^{\mathrm{ref}} \in \mathbb{R}^{m}$ is the reference point and $\mathcal{F}^{(G)}$
denotes the set of objective vectors associated with the population $\mathcal{U}^{(G)}$.
\nomenclature{$\mathrm{HV}(\cdot,\cdot)$}{Hypervolume indicator of an objective set with respect to a reference point}
\nomenclature{$\mathrm{HV}^{(G)}$}{Hypervolume indicator at evolutionary generation $G$}
\nomenclature{$\mathcal{F}^{(G)}$}{Objective vectors of the evolutionary population at generation $G$}
\nomenclature{$\mathbf{z}^{\mathrm{ref}}$}{Reference point for hypervolume computation}

Equivalently, when only the non-dominated objective vectors available after generation $G$ are
retained, denoted $\mathcal{F}_{\mathrm{EA}}^{(G)}$, the corresponding hypervolume is
\begin{equation}
    \begin{gathered}
        \mathrm{HV}_{\mathrm{EA}}^{(G)}
        =
        \mathrm{HV}
        \left(
            \mathcal{F}_{\mathrm{EA}}^{(G)},
            \mathbf{z}^{\mathrm{ref}}
        \right), \\[4pt]
        \mathcal{F}_{\mathrm{EA}}^{(G)} \subseteq \mathcal{F}^{(G)} \subset \mathbb{R}^{m}
        \;\text{(non-dominated subset)},\quad
        1 \leq \bigl|\mathcal{F}_{\mathrm{EA}}^{(G)}\bigr| \leq N_{\mathrm{pop}},\\[2pt]
        \mathrm{HV}_{\mathrm{EA}}^{(G)} \in \mathbb{R}_{\geq 0};
        \quad
        \mathrm{HV}_{\mathrm{EA}}^{(G+1)} \geq \mathrm{HV}_{\mathrm{EA}}^{(G)}
        \;\text{only if } \mathcal{F}_{\mathrm{EA}}^{(G)}
        \text{ is a cumulative archive}
    \end{gathered}
\end{equation}

For Bayesian optimization, the hypervolume can be calculated after every completed trial or after
every completed candidate batch. Let $\mathcal{F}_{\mathrm{BO}}^{(t)}$ denote the non-dominated
objective vectors obtained from the first $t$ completed Bayesian trials. The corresponding
hypervolume is
\begin{equation}
    \begin{gathered}
        \mathrm{HV}_{\mathrm{BO}}^{(t)}
        =
        \mathrm{HV}
        \left(
            \mathcal{F}_{\mathrm{BO}}^{(t)},
            \mathbf{z}^{\mathrm{ref}}
        \right), \\[4pt]
        \mathcal{F}_{\mathrm{BO}}^{(t)} \subset \mathbb{R}^{m},\quad
        \bigl|\mathcal{F}_{\mathrm{BO}}^{(t)}\bigr| \leq t,\quad
        t \in \mathbb{N}_0,\quad
        \mathrm{HV}_{\mathrm{BO}}^{(0)} = 0,\\[2pt]
        \mathrm{HV}_{\mathrm{BO}}^{(t+1)} \geq \mathrm{HV}_{\mathrm{BO}}^{(t)}
        \;\text{(cumulative archive by construction)}
    \end{gathered}
\end{equation}
\nomenclature{$\mathcal{F}_{\mathrm{EA}}^{(G)}$}{Non-dominated objective set after evolutionary generation $G$}
\nomenclature{$\mathcal{F}_{\mathrm{BO}}^{(t)}$}{Non-dominated objective set after $t$ completed Bayesian evaluations}
\nomenclature{$\mathrm{HV}_{\mathrm{EA}}^{(G)}$}{Evolutionary hypervolume after generation $G$}
\nomenclature{$\mathrm{HV}_{\mathrm{BO}}^{(t)}$}{Bayesian hypervolume after $t$ completed evaluations}

Under the maximisation convention, the fixed reference point $\mathbf{z}^{\mathrm{ref}}$ must be
worse than the objective values of interest. The same reference point and objective scaling must be
used for all compared methods. An increasing hypervolume generally indicates progress toward a
better approximation of the Pareto front, improved coverage of the objective space, or both.
Monotonicity is guaranteed only if the indicator is computed from a cumulative non-dominated
archive; the hypervolume of the current population may fluctuate between generations.

Evolutionary generations and Bayesian batches can contain different numbers of candidate
evaluations. Therefore, the primary comparison is performed using the hypervolume as a function of
the cumulative number of objective evaluations,
\begin{equation}
    \begin{gathered}
        \mathrm{HV} = \mathrm{HV}(E), \\[4pt]
        \mathrm{HV}(\cdot) : \{0,1,\ldots,E_{\max}^{\bullet}\} \rightarrow \mathbb{R}_{\geq 0},
        \qquad
        \bullet \in \{\mathrm{EA},\mathrm{BO}\},\\[2pt]
        E \in \mathbb{N}_0,\quad
        \mathrm{HV}(0) = 0
    \end{gathered}
\end{equation}
where $E$ is the number of completed external API evaluations.
\nomenclature{$E$}{Cumulative number of completed external objective evaluations}


\subsection{Optional Hyperparameter Optimization}
\label{sec:two_finger_hyperparameter_optimization}

For the evolutionary algorithms, the framework optionally supports automated tuning of algorithm
hyperparameters, such as crossover probability, mutation rate, mutation eta, crossover eta,
selection eta, etc., using a meta-optimization approach. A \texttt{HyperparameterProblem} is
constructed around the selected algorithm and evaluated over multiple random seeds. The negative
mean hypervolume across the runs is used as the meta-objective:
\begin{equation}
    \begin{gathered}
        \min_{\boldsymbol{\theta}\in\Theta}
        \left\{
            -\frac{1}{|\mathcal{Z}|}
            \sum_{\sigma \in \mathcal{Z}}
            \mathrm{HV}\!\left(\mathcal{F}_{\sigma}(\boldsymbol{\theta}),\mathbf{z}^{\mathrm{ref}}\right)
        \right\}, \\[4pt]
        \Theta = \Theta_{\mathrm{cont}} \times \Theta_{\mathrm{int}} \times \Theta_{\mathrm{cat}}
        \;\text{(mixed hyperparameter space)},\quad
        \mathcal{Z} \subset \mathbb{N},\;
        |\mathcal{Z}| \in \mathbb{N},\\[2pt]
        \mathcal{F}_{\sigma}(\boldsymbol{\theta}) \subset \mathbb{R}^{m}
        \;\text{finite},\quad
        \text{objective value} \in \left[-\mathrm{HV}_{\max},\,0\right]
    \end{gathered}
\end{equation}
\nomenclature{$\boldsymbol{\theta}$}{Algorithm hyperparameters of the evolutionary algorithm}
\nomenclature{$\mathcal{Z}$}{Set of random seeds used for meta-optimization}
\nomenclature{$\sigma$}{Random-seed index}
\nomenclature{$\mathcal{F}_{\sigma}$}{Final objective vectors obtained with random seed $\sigma$}

The best hyperparameter configuration is subsequently applied to the main evolutionary optimization
run.

The current Bayesian implementation uses the generation strategy, surrogate configuration, and
acquisition-function configuration selected by Ax. A separate meta-optimization of Bayesian model
hyperparameters is therefore not currently performed. The user specifies the candidate batch size
and the optimization budget, while Ax manages the transition from initialization to model-based
optimization.

\subsection{Output}
\label{sec:common_output}

Because all objectives are maximised in their natural form, the optimization produces the set of
Pareto-optimal solutions
\begin{equation}
    \mathcal{P}^*
    =
    \left\{
        \mathbf{q} \in \mathcal{Q}
        \;\middle|\;
        \nexists\; \mathbf{q}' \in \mathcal{Q}
        \text{ such that }
        \mathbf{f}(\mathbf{q}') \geq \mathbf{f}(\mathbf{q})
        \text{ component-wise and }
        \mathbf{f}(\mathbf{q}') \neq \mathbf{f}(\mathbf{q})
    \right\}.
\end{equation}
\nomenclature{$\mathcal{P}^*$}{Set of Pareto-optimal solutions in the decision space}

Equivalently, $\mathbf{q}'\succ\mathbf{q}$ denotes that $\mathbf{q}'$ dominates $\mathbf{q}$, i.e.
\begin{equation}
    \mathcal{P}^{*}
    =
    \left\{
        \mathbf{q}\in\mathcal{Q}
        \;\middle|\;
        \nexists\,\mathbf{q}'\in\mathcal{Q}:
        \mathbf{q}' \succ \mathbf{q}
    \right\},
\end{equation}
with
\begin{equation}
    f_i(\mathbf{q}')\geq f_i(\mathbf{q})
    \quad
    \forall\, i = 1,\ldots,m
\end{equation}
and
\begin{equation}
    f_j(\mathbf{q}')>f_j(\mathbf{q})
    \quad
    \text{for at least one objective } j .
\end{equation}

The final output of an optimization run includes:
\begin{itemize}
    \item The optimization configuration used to carry out the optimization from
          Section~\ref{sec:optimization_framework_sc}
    \item The optimization algorithm hyperparameters
    \item The Pareto-optimal set $\mathcal{P}^*$ in the decision space
    \item The corresponding objective values $\mathbf{f}(\mathbf{q})$
    \item The hypervolume progression
    \item Total optimization wall-clock time
\end{itemize}
\nomenclature{$G_{\mathrm{final}}$}{Final evolutionary generation reached by the optimization}

The problem-specific contents and interpretations of the returned decision vectors are given in
Section~\ref{sec:multi_sc_output} and Section~\ref{sec:two_finger_output}.

% ============================================================================
\section{Implementation of Multi-Suction Cup Gripper Optimization}
\label{sec:multi_sc_gripper_optimization}

This section describes the problem-specific implementation for a multi-suction cup gripper
consisting of three suction cups arranged in a triangular configuration. The common framework
components are described in Section~\ref{sec:gripper_optimization_framework}.

\subsection{Reduced Parameterization of Multi-Suction Cup Gripper}

The gripper consists of three suction cups whose positions must ultimately be represented in
Cartesian space in order to perform grasp generation, collision checking, and objective computation.
However, the Cartesian suction cup coordinates are not used directly as optimization variables.
Instead, the layout is parameterized in the reduced form introduced in
Section~\ref{subsec:methodology_decision_variables} in order to constrain the search space and
encode geometric structure.

All three suction cups are constrained to lie on a common circle in the gripper plane. Their
configuration is therefore defined by one shared radius $r$ and three angular parameters
$\alpha_1,\alpha_2,\alpha_3$:
\begin{equation}
    \mathbf{q}_{\mathrm{SC}} = [r,\alpha_1,\alpha_2,\alpha_3].
\end{equation}
For brevity, $\mathbf{q}$ denotes the suction-cup decision vector $\mathbf{q}_{\mathrm{SC}}$
throughout this section.

The angular parameters determine the positions of the suction cups on the circle, while the common
radius determines their distance from the gripper Tool Center Point (TCP). For a fixed gripper
height $z_0 = 0.326$~m, the Cartesian coordinates of the suction cups are obtained through the
polar-to-Cartesian transformation
\begin{equation}
    \mathbf{p}_i =
    \begin{bmatrix}
        r\cos(\alpha_i) \\
        r\sin(\alpha_i) \\
        z_0
    \end{bmatrix},
    \qquad i=1,2,3.
\end{equation}
\nomenclature{$z_0$}{Fixed height coordinate of the suction cups in the gripper frame}

The resulting Cartesian gripper configuration is written as
\begin{equation}
    \mathbf{P} = [\mathbf{p}_1,\mathbf{p}_2,\mathbf{p}_3].
\end{equation}
\nomenclature{$\mathbf{P}$}{Cartesian suction cup configuration derived from the reduced parameterization}

Thus, the optimization is carried out in the reduced four-dimensional parameter space of
$\mathbf{q}$, while the grasp evaluation operates on the derived Cartesian configuration
$\mathbf{P}$.

This formulation reduces the dimensionality of the search space from six planar degrees of freedom,
\begin{equation}
    \left[(p_{1,x},p_{1,y}),(p_{2,x},p_{2,y}),(p_{3,x},p_{3,y})\right] \in \mathbb{R}^6,
\end{equation}
to four variables,
\begin{equation}
    [r,\alpha_1,\alpha_2,\alpha_3].
\end{equation}
\nomenclature{$p_{i,x},p_{i,y}$}{Planar coordinates of suction cup $i$ in the gripper frame}

The reduction improves search efficiency and excludes many geometrically redundant or undesirable
layouts while preserving sufficient flexibility for diverse triangular suction cup arrangements.

\paragraph{Assumption}
A triangular configuration of three suction cups is assumed to yield a stable
grasp. This is motivated by the fact that three spatially
distributed contact points form a statically stable support polygon, enabling effective load
distribution and resistance to external disturbances.

\subsection{Decision Variables and Bounds for Multi-Suction Cup Gripper}
\label{sec:multi_sc_decision_variables}

The optimization is performed over the reduced decision vector
\begin{equation}
    \mathbf{q} = [r,\alpha_1,\alpha_2,\alpha_3].
\end{equation}

All variables are treated as continuous. The angular variables are bounded by
\begin{equation}
    \alpha_i \in [0,2\pi), \qquad i=1,2,3,
\end{equation}
and the common radius is bounded by
\begin{equation}
    r \in [r_{\min},r_{\max}].
\end{equation}
\nomenclature{$r_{\min},r_{\max}$}{Lower and upper bound on the common suction cup radius}

The radial bounds are determined separately for each workpiece. The workpiece is oriented such that
its longest dimension lies in the horizontal plane, the horizontal bounding box is centred at
\((x,y)=(0,0)\), and the base of the workpiece is placed at \(z=0\). This establishes a consistent
coordinate system for defining the gripper-design bounds. Let \(X_{\min}\), \(X_{\max}\),
\(Y_{\min}\), and \(Y_{\max}\) denote the coordinates of the bounding box of a workpiece in the
\(x\)- and \(y\)-directions. The maximum radial distance is defined by the distance from the origin
to the nearest bounding-box boundary:
\begin{equation}
    r_{\max}
    =
    \min
    \left\{
        |X_{\min}|,\,
        |X_{\max}|,\,
        |Y_{\min}|,\,
        |Y_{\max}|
    \right\}.
\end{equation}
\nomenclature{$X_{\min},X_{\max},Y_{\min},Y_{\max}$}{Horizontal bounding-box coordinates of the workpiece}

The lower radial bound is derived from the suction-cup radius \(r_{\mathrm{cup}}\) and is defined as
\begin{equation}
    r_{\min}
    =
    \frac{2\sqrt{3}}{3}\,
    r_{\mathrm{cup}} .
\end{equation}
\nomenclature{$r_{\mathrm{cup}}$}{Radius of a single suction cup}

This lower bound corresponds to the circumradius of an equilateral triangle with side length
\(2r_{\mathrm{cup}}\), i.e.\ the configuration in which neighbouring cups just touch.

All optimization methods generate candidate values within these bounds. In the evolutionary
algorithms, the initial population is obtained by random sampling and subsequently modified through
the operators of the selected algorithm. In Bayesian optimization, Ax first evaluates the center of
the search space and then uses Sobol sampling to generate an initial space-filling design. After
sufficient initialization data have been collected, Ax transitions to model-based Bayesian
optimization.

\subsection{Objective Function Evaluation of Multi-Suction Cup Gripper via API}
\label{sec:multi_sc_api_evaluation}

The evaluation of objective values is decoupled from the optimization loop and delegated to an
external evaluation service. For each candidate solution $\mathbf{q}$, the corresponding Cartesian
suction cup coordinates are first computed by transforming the radial and angular parameters into 3D
positions. The resulting Cartesian gripper configuration is then passed to the evaluation endpoint,
which executes the grasp evaluation procedure and returns the objective values.

This process can be written as
\begin{equation}
    \mathbf{P} = \Phi(\mathbf{q}),
    \qquad
    \mathbf{f}(\mathbf{q})
    =
    \texttt{API\_call}(\mathbf{P})
    =
    \left[D,\;\hat{n}_{\mathrm{g}},\;\hat{A}\right],
\end{equation}
where $\Phi(\cdot)$ denotes the transformation from the reduced parameterization to Cartesian
suction cup coordinates, and the three objective components $D$, $\hat{n}_{\mathrm{g}}$ and
$\hat{A}$ are defined in Section~\ref{sec:grasp_eval_sc}. Since every API call evaluates exactly one
fixed configuration $\mathbf{P}$, the argument is omitted from the individual quality measures
wherever it is unambiguous.
\nomenclature{$\Phi(\cdot)$}{Transformation from reduced decision variables to Cartesian suction cup coordinates}

In the implementation, this transformation is carried out programmatically in the evaluation
pipeline before grasp generation and objective computation. For a candidate radius and list of
angles, the suction cup coordinates are constructed as
\begin{equation}
    \mathbf{p}_i = \left(r\cos(\alpha_i),\; r\sin(\alpha_i),\; z_0\right),
    \qquad i=1,2,3.
\end{equation}

For every candidate, the evaluation request performs the following steps:
\begin{enumerate}
    \item Send the candidate decision vector to the external evaluation API.
    \item Transform $[r,\alpha_1,\alpha_2,\alpha_3]$ into the Cartesian configuration
          $\mathbf{P}=\Phi(\mathbf{q})$.
    \item Generate and filter the suction grasp candidates by seal formation and collision check.
    \item Cluster the valid grasps and compute $D$, $\hat{n}_{\mathrm{g}}$ and $\hat{A}$.
    \item Return the objective vector $\mathbf{f}(\mathbf{q})$ to the selected optimization
          framework.
\end{enumerate}

This architecture enables:
\begin{itemize}
    \item Separation of concerns between optimization logic and domain-specific evaluation
    \item Scalability through independent deployment of the evaluation service
    \item Reusability of the optimization framework for different problem domains
\end{itemize}

\subsubsection{Grasp Evaluation}
\label{sec:grasp_eval_sc}

The evaluation service receives a Cartesian gripper configuration
$\mathbf{P} = [\mathbf{p}_1,\mathbf{p}_2,\mathbf{p}_3]$ and computes the associated grasp quality
measures through the following procedure.

\begin{enumerate}
    \item \textbf{Candidate Grasp Generation and Feasibility Filtering}

    Suction grasp candidates are generated using a method developed by Fraunhofer IPA, following the
    automatic suction grasp generation procedure described by~\cite{Kleeberger}. Points are
    uniformly sampled on the surface of the workpiece, and each sampled point defines a candidate
    grasp through its position and local surface normal. The local surface normal determines the
    approach direction of the suction cup.

    In the original method, each suction grasp candidate is tested for feasible seal formation using
    a quasi-static spring model of a suction cup, as introduced in~\cite{SC_quasi_static}. If the
    grasp point is considered feasible, the candidate is further tested with a gripper collider
    model for collisions between the gripper and the object.

    In this thesis, the method was extended to a gripper with three suction cups. For each candidate
    gripper pose, the seal-formation test is performed for all three suction cup contact positions
    $\mathbf{p}_1$, $\mathbf{p}_2$, and $\mathbf{p}_3$. A candidate is considered feasible only if
    all three suction cups form a sufficient seal. Each feasible candidate is subsequently subjected
    to a geometric collision check: the gripper geometry is represented by a sampled point set,
    transformed into the workpiece frame by the candidate pose, and tested for proximity against the
    workpiece point cloud. A candidate is rejected if any transformed gripper point lies closer than
    the clearance threshold $\delta_{\mathrm{clr}} = 0.5\,\text{mm}$ to any workpiece point. The
    candidates that pass both the seal-formation and the collision checks form the set of valid
    grasps $\mathcal{H}_{\mathrm{valid}}$, which is the input to the subsequent clustering step.
    \nomenclature{$\delta_{\mathrm{clr}}$}{Clearance threshold of the gripper--workpiece collision check}
    \nomenclature{$\mathcal{H}_{\mathrm{valid}}$}{Set of valid grasps after seal-formation and collision filtering}

    \item \textbf{Clustering of Grasps}

    The set $\mathcal{H}_{\mathrm{valid}}$ typically contains many grasps that are nearly identical
    in position and orientation. If the number of valid grasps
    $n_{\mathrm{v}} = |\mathcal{H}_{\mathrm{valid}}|$ exceeds the clustering budget
    $K_{\mathrm{clst}} = 15$, the grasps are clustered using the greedy farthest-point sampling
    (FPS) algorithm of~\cite{gonzalez1985clustering}, which iteratively selects the grasp with the
    largest minimum distance to the already selected representatives; otherwise all valid grasps are
    retained.
    \nomenclature{$n_{\mathrm{v}}$}{Number of valid grasps returned by the feasibility filtering}
    \nomenclature{$K_{\mathrm{clst}}$}{Clustering budget, i.e.\ maximum number of grasps retained by clustering}

    Since grasp poses combine a translational and a rotational component, the sampling operates on a
    combined pose metric in which the Euclidean distance between the grasp positions
    $\mathbf{x}_i \in \mathbb{R}^3$ is normalised by the diagonal extent of the axis-aligned
    bounding box of all candidates, and the difference between the orientations
    $\mathbf{R}_i \in SO(3)$ is measured as the geodesic angle normalised by $\pi$~\cite{Huynh2009}.
    With the resulting normalised terms $\hat{\rho}_{\mathrm{pos}}$ and $\hat{\rho}_{\mathrm{ori}}$
    and an orientation weight $w \in [0,1]$,
    \begin{equation}
        \label{eq:combined-metric}
        \rho_w(i,j) =
        (1 - w)\,\hat{\rho}_{\mathrm{pos}}(i,j) + w\,\hat{\rho}_{\mathrm{ori}}(i,j),
    \end{equation}
    so that $w = 0$ corresponds to purely position-based and $w = 1$ to purely orientation-based
    sampling; here $w = 0.35$ is used. All distances are evaluated on the GPU, which keeps the
    selection tractable for large candidate sets.

    The final number of representatives $K^{*}$ is not predefined but obtained from an elbow
    criterion on the covering radius recorded during sampling: selection stops as soon as the
    relative coverage gain of an additional representative falls below $\tau = 0.05$, bounded from
    above by $K_{\mathrm{clst}}$. The resulting subset
    $\mathcal{H}_{\mathrm{clst}} \subseteq \mathcal{H}_{\mathrm{valid}}$ with
    $|\mathcal{H}_{\mathrm{clst}}| = K^{*}$ is passed on to the diversity evaluation.
    \nomenclature{$\mathbf{R}_i$}{Orientation matrix of grasp $i$}
    \nomenclature{$\mathbf{x}_i$}{Position of grasp $i$}
    \nomenclature{$\rho_w(\cdot,\cdot)$}{Combined pose distance with orientation weight $w$}
    \nomenclature{$w$}{Orientation weight in the combined pose distance}
    \nomenclature{$K^{*}$}{Final number of clustered grasps selected by the elbow criterion}

    \item \textbf{Diversity Metric Calculation}

    A diversity score $D$ is computed to quantify how broadly the resulting valid grasp set is
    distributed over the object, both in position and in orientation.

    Let the set used for diversity evaluation contain $n_{\mathrm{g}}$ grasps. If clustering is
    performed, then
    \begin{equation}
        n_{\mathrm{g}} = |\mathcal{H}_{\mathrm{clst}}|,
    \end{equation}
    otherwise
    \begin{equation}
        n_{\mathrm{g}} = |\mathcal{H}_{\mathrm{valid}}| = n_{\mathrm{v}} .
    \end{equation}
    For $n_{\mathrm{g}} < 2$ no pairwise comparison is defined; in this case the corresponding
    component is set to zero.

    The evaluated positions are the projections of the grasp points onto the mesh surface,
    \begin{equation}
        \mathbf{y}_i = \arg\min_{\boldsymbol{\xi} \in \partial \mathcal{W}}
        \|\boldsymbol{\xi} - \mathbf{x}_i\|_2 ,
        \qquad i = 1,\dots,n_{\mathrm{g}},
    \end{equation}
    where $\partial \mathcal{W}$ denotes the surface of the workpiece mesh and $\mathbf{x}_i$ the
    original grasp point. The associated orientation matrices are
    $\{\mathbf{R}_1,\dots,\mathbf{R}_{n_{\mathrm{g}}}\}$, with $\mathbf{R}_i \in SO(3)$.
    \nomenclature{$\mathcal{W}$}{Workpiece mesh, with surface $\partial\mathcal{W}$}

    To normalise the spatial terms with respect to object size, a characteristic object scale is
    defined as the Euclidean norm of the axis-aligned bounding-box extents,
    \begin{equation}
        s = \|\mathbf{e}\|_2 ,
    \end{equation}
    where $\mathbf{e} \in \mathbb{R}^3$ contains the bounding-box side lengths.

    The spatial diversity is assessed by two complementary quantities. The first is the normalised
    mean pairwise Euclidean distance between the grasp positions,
    \begin{equation}
        \begin{gathered}
            D_{\mathrm{spread}} =
            \frac{1}{s} \cdot
            \frac{1}{n_{\mathrm{g}}(n_{\mathrm{g}}-1)}
            \sum_{\substack{i,j=1 \\ i \neq j}}^{n_{\mathrm{g}}}
            \|\mathbf{y}_i - \mathbf{y}_j\|_2 , \\[4pt]
            \mathbf{y}_i \in \partial\mathcal{W} \subset \mathbb{R}^{3},\quad
            s = \|\mathbf{e}\|_2 \in \mathbb{R}_{>0},\quad
            \mathbf{e} \in \mathbb{R}_{>0}^{3},\quad
            n_{\mathrm{g}} \geq 2,\\[2pt]
            \|\mathbf{y}_i-\mathbf{y}_j\|_2 \leq s
            \;\Longrightarrow\;
            D_{\mathrm{spread}} \in [0,1],\quad
            D_{\mathrm{spread}} := 0 \text{ for } n_{\mathrm{g}} < 2
        \end{gathered}
    \end{equation}
    Because a large mean distance can also be achieved by grasps concentrated on two opposite faces,
    a second term evaluates the angular coverage of the grasp positions as seen from the object
    centroid $\mathbf{o}$. With the normalised view directions
    \begin{equation}
        \begin{gathered}
            \hat{\mathbf{u}}_i =
            \frac{\mathbf{y}_i - \mathbf{o}}{\|\mathbf{y}_i - \mathbf{o}\|_2 + \varepsilon}, \\[4pt]
            \mathbf{y}_i,\mathbf{o} \in \mathbb{R}^{3},\quad
            \varepsilon \in \mathbb{R}_{>0},\; \varepsilon \ll 1,\quad
            \hat{\mathbf{u}}_i \in \mathbb{R}^{3},\quad
            \|\hat{\mathbf{u}}_i\|_2 \in [0,1)
        \end{gathered}
    \end{equation}
    where $\varepsilon$ is a small constant guarding against division by zero, the pairwise cosine
    distance
    \begin{equation}
        \rho_{\cos}(i,j) = 1 - \hat{\mathbf{u}}_i^\top \hat{\mathbf{u}}_j \in [0,2]
    \end{equation}
    vanishes for identical directions and attains its maximum for opposing directions. Its
    normalised mean defines the coverage term
    \begin{equation}
        D_{\mathrm{cov}} =
        \mathrm{clip}\!\left(
            \frac{1}{2} \cdot
            \frac{1}{n_{\mathrm{g}}(n_{\mathrm{g}}-1)}
            \sum_{\substack{i,j=1 \\ i \neq j}}^{n_{\mathrm{g}}}
            \rho_{\cos}(i,j),
            \; 0, \; 1
        \right).
    \end{equation}
    Both contributions are combined using the angular-coverage weight $\beta_{\mathrm{cov}} = 0.5$:
\begin{equation}
    \begin{gathered}
        D_{\mathrm{pos}} =
        (1-\beta_{\mathrm{cov}})\,D_{\mathrm{spread}}
        +
        \beta_{\mathrm{cov}}\,D_{\mathrm{cov}}, \\[4pt]
        \beta_{\mathrm{cov}} = 0.5 \in [0,1],\quad
        D_{\mathrm{spread}},D_{\mathrm{cov}} \in [0,1]
        \;\Longrightarrow\;
        D_{\mathrm{pos}} \in [0,1]
    \end{gathered}
\end{equation}
    so that a high spatial score requires the grasps to be both mutually distant and distributed
    around the object rather than clustered on a single side.

    The orientation component is computed from the mean pairwise geodesic distance between rotation
    matrices. For two grasp orientations $\mathbf{R}_i, \mathbf{R}_j \in SO(3)$, the relative
    rotation is
    \begin{equation}
        \Delta \mathbf{R}_{ij} = \mathbf{R}_i^\top \mathbf{R}_j
    \end{equation}
    and the associated geodesic distance is
\begin{equation}
    \begin{gathered}
        \rho_{\mathrm{geo}}(\mathbf{R}_i,\mathbf{R}_j)
        = \cos^{-1}\!\left(
            \frac{\mathrm{tr}(\Delta \mathbf{R}_{ij}) - 1}{2}
        \right), \\[4pt]
        \mathbf{R}_i,\mathbf{R}_j \in \mathrm{SO}(3),\quad
        \mathrm{tr}(\Delta\mathbf{R}_{ij}) \in [-1,3]
        \;\Longrightarrow\;
        \tfrac{1}{2}\!\left(\mathrm{tr}(\Delta\mathbf{R}_{ij})-1\right) \in [-1,1],\\[2pt]
        \rho_{\mathrm{geo}} : \mathrm{SO}(3)\times\mathrm{SO}(3) \rightarrow [0,\pi],\quad
        \rho_{\mathrm{geo}}(\mathbf{R},\mathbf{R}) = 0
    \end{gathered}
\end{equation}
    where the argument of $\cos^{-1}$ is clipped to $[-1,1]$ for numerical stability. The normalised
    orientation diversity follows as
    \begin{equation}
        D_{\mathrm{rot}} =
        \frac{1}{\pi} \cdot
        \frac{2}{n_{\mathrm{g}}(n_{\mathrm{g}}-1)}
        \sum_{1 \leq i < j \leq n_{\mathrm{g}}}
        \rho_{\mathrm{geo}}(\mathbf{R}_i,\mathbf{R}_j).
    \end{equation}

    The final diversity score weights the spatial component more strongly than the orientation
    component, since the orientation of suction grasps is largely determined by the local surface
    normal and therefore varies less independently. With $\lambda = 0.65$,
\begin{equation}
    \begin{gathered}
        D =
        \mathrm{clip}\!\left(
            \lambda\, D_{\mathrm{pos}} +
            (1-\lambda)\, D_{\mathrm{rot}},
            \; 0, \; 1
        \right), \\[4pt]
        \lambda = 0.65 \in [0,1],\quad
        D_{\mathrm{pos}},D_{\mathrm{rot}} \in [0,1]
        \;\Longrightarrow\;
        \lambda D_{\mathrm{pos}} + (1-\lambda)D_{\mathrm{rot}} \in [0,1],\\[2pt]
        D \in [0,1]
        \;\text{(the clip only compensates floating-point round-off)}
    \end{gathered}
\end{equation}
    where the final clipping prevents values marginally exceeding unity due to floating-point
    rounding. For large grasp sets the orientation term saturates around
    $D_{\mathrm{rot}} \approx 0.7$, as uniformly distributed rotations do not reach the theoretical
    maximum of the normalised geodesic mean.
    \nomenclature{$D(\cdot)$}{Diversity metric of the grasp distribution of a gripper layout}
    \nomenclature{$D_{\mathrm{pos}}$}{Spatial component of the diversity metric}
    \nomenclature{$D_{\mathrm{spread}}$}{Normalised mean pairwise distance between grasp positions}
    \nomenclature{$D_{\mathrm{cov}}$}{Angular coverage of grasp positions around the object centroid}
    \nomenclature{$D_{\mathrm{rot}}$}{Orientation component of the diversity metric}
    \nomenclature{$\rho_{\mathrm{geo}}(\cdot,\cdot)$}{Geodesic distance between two grasp orientations}
    \nomenclature{$\rho_{\cos}(\cdot,\cdot)$}{Cosine distance between the centroid-referenced directions of two grasps}
    \nomenclature{$\Delta\mathbf{R}_{ij}$}{Relative rotation between the orientations of grasps $i$ and $j$}
    \nomenclature{$\hat{\mathbf{u}}_i$}{Normalised direction from the object centroid to grasp position $i$}
    \nomenclature{$\mathbf{o}$}{Centroid of the workpiece mesh}
    \nomenclature{$\varepsilon$}{Small numerical constant preventing division by zero}
    \nomenclature{$\beta_{\mathrm{cov}}$}{Weight of the angular coverage term within the spatial diversity component}
    \nomenclature{$\lambda$}{Weight of the spatial component in the total diversity score}
    \nomenclature{$s$}{Characteristic object scale derived from the mesh bounding box}
    \nomenclature{$\mathbf{e}$}{Axis-aligned bounding-box side lengths of the workpiece}
    \nomenclature{$\mathbf{y}_i$}{Surface-projected position of grasp $i$ used in diversity evaluation}
    \nomenclature{$n_{\mathrm{g}}$}{Number of grasps used for diversity evaluation}

    \item \textbf{Normalised Grasp Count}

    Besides the spatial distribution of the grasps, their absolute number is relevant as well, since
    a layout that offers more valid grasps is more robust with respect to object pose and
    surrounding scene constraints. The evaluated grasp count is taken from the same set that is used
    for the diversity metric, i.e.
    \begin{equation}
        n_{\mathrm{g}} =
        \begin{cases}
            |\mathcal{H}_{\mathrm{clst}}|,  & \text{if clustering is performed},\\[2pt]
            |\mathcal{H}_{\mathrm{valid}}|, & \text{otherwise}.
        \end{cases}
    \end{equation}
    To keep this quantity comparable with the remaining objectives, it is expressed relative to the
    clustering budget $K_{\mathrm{clst}}$ (\texttt{max\_cluster}), which is the maximum number of
    grasps retained by the clustering stage,
    \begin{equation}
        \hat{n}_{\mathrm{g}} = \frac{n_{\mathrm{g}}}{K_{\mathrm{clst}}} \in [0,1].
    \end{equation}
    Since clustering is only applied when more than $K_{\mathrm{clst}}$ valid grasps exist, the
    condition $n_{\mathrm{g}} \leq K_{\mathrm{clst}}$ holds in both cases and the ratio is bounded
    by unity without further clipping. A value of one indicates that the layout produced at least as
    many valid grasps as the clustering budget permits, whereas a value of zero corresponds to a
    layout for which no valid grasp was found.
    \nomenclature{$\hat{n}_{\mathrm{g}}(\cdot)$}{Normalised grasp count of a gripper layout, i.e.\ the number of evaluated grasps divided by the clustering budget}

    \item \textbf{Triangle Area Calculation}

    The area of the triangle formed by the three suction cups is computed from their Cartesian
    coordinates. With the planar coordinates $(p_{i,x},p_{i,y})$ of cup $i$,
\begin{equation}
    \begin{gathered}
        A =
        \frac{1}{2}
        \left|
            (p_{2,x} - p_{1,x})(p_{3,y} - p_{1,y}) - (p_{3,x} - p_{1,x})(p_{2,y} - p_{1,y})
        \right|, \\[4pt]
        \mathbf{P} \in \mathcal{P}_{\mathcal{W}},\quad
        p_{i,x},p_{i,y} \in \left[-r_{\max},\,r_{\max}\right],\quad
        A : \mathcal{P}_{\mathcal{W}} \rightarrow \left[0,\,A_{\max}\right],\\[2pt]
        A_{\max} = \tfrac{3\sqrt{3}}{4}r_{\max}^{2},\quad
        A = 0 \iff \text{the three cups are collinear}
    \end{gathered}
\end{equation}

    To make this objective comparable in magnitude with the remaining objectives, the area is
    normalised by the largest triangle area that can be realised within the admissible mounting
    region. The suction-cup reference points are restricted to a circle of radius $r_{\max}$, and
    the maximum-area triangle inscribed in a circle is the equilateral one, so that
    \begin{equation}
        A_{\max} = \frac{3\sqrt{3}}{4}\, r_{\max}^{2}.
    \end{equation}
    The radius $r_{\max}$ is not a global constant but is determined per workpiece as the largest
    circle that still fits inside the admissible mounting bounds. Consequently $\hat{A}$ expresses
    the achieved area relative to what is geometrically possible for that specific workpiece.

    The normalised area objective is therefore
    \begin{equation}
        \hat{A} = \frac{A}{A_{\max}} \in [0,1].
    \end{equation}

    This objective favours gripper layouts with a wider support polygon, which is associated with
    the assumption of improved grasp stability.
    \nomenclature{$A(\cdot)$}{Area of the triangle formed by the three suction cups of a gripper layout}
    \nomenclature{$\hat{A}(\cdot)$}{Normalised triangle-area objective of a gripper layout}
    \nomenclature{$A_{\max}$}{Maximum triangle area attainable within the admissible mounting circle}
\end{enumerate}

\paragraph{Continuity considerations}
Clustering is only applied when the number of valid grasps exceeds the clustering budget
$K_{\mathrm{clst}}$ (\texttt{max\_cluster}). If fewer valid grasps are available, clustering is
skipped and the diversity metric is evaluated directly on the valid grasp set. Therefore, the
omission of clustering does not by itself introduce a discontinuity in the objective function.

A discontinuity may arise only in the degenerate case where fewer than two valid grasps remain after
collision filtering. Since both the spatial and orientation components of the diversity metric are
based on pairwise distances, these terms cannot be evaluated meaningfully for $n_{\mathrm{v}} < 2$.
In that case, the corresponding terms are set to zero. As a result, small changes in the gripper
configuration may cause an abrupt transition between a non-zero diversity score and a degenerate
zero-valued diversity contribution.

To reduce the likelihood of this case, the number of initially sampled surface points was increased
from 100 to 1000. The denser surface sampling produces a larger set of candidate grasps and makes it
much less likely that fewer than two valid grasps remain after filtering.

\subsection{Objective Evaluation of Multi-Suction Cup Gripper}
\label{sec:multi_sc_objective_evaluation}

After grasp evaluation, each candidate layout is assigned the natural objective vector
\begin{equation}
    \mathbf{f}(\mathbf{q})
    =
    \left[
        D(\mathbf{P}),\;
        \hat{n}_{\mathrm{g}}(\mathbf{P}),\;
        \hat{A}(\mathbf{P})
    \right],
    \qquad
    \mathbf{P} = \Phi(\mathbf{q}),
\end{equation}
where the ordering corresponds to the return order of the evaluation pipeline, so that $m = 3$ for
this problem. In the remainder of the section the three components are written without arguments
whenever the underlying layout is clear from the context.

All three components are maximised:
\begin{itemize}
    \item $D$ -- the diversity score, where larger values indicate that the valid grasps are more
          broadly distributed over the object in position and orientation.
    \item $\hat{n}_{\mathrm{g}}$ -- the normalised grasp count, where larger values indicate that
          the layout can grasp the object in more distinct ways and therefore offers greater
          robustness with respect to object pose and scene constraints.
    \item $\hat{A}$ -- the normalised triangle area spanned by the three suction cups, where larger
          values correspond to a wider support polygon.
\end{itemize}

Following the convention of Section~\ref{sec:objective_direction_conventions}, the evolutionary
algorithms receive
\begin{equation}
    \mathbf{f}_{\mathrm{EA}}(\mathbf{q})
    =
    -\,\mathbf{f}(\mathbf{q})
    =
    \left[-D,\;-\hat{n}_{\mathrm{g}},\;-\hat{A}\right],
\end{equation}
whereas the Bayesian implementation uses the unchanged observations
\begin{equation}
    \mathbf{f}_{\mathrm{BO}}(\mathbf{q})
    =
    \mathbf{f}(\mathbf{q})
    =
    \left[D,\;\hat{n}_{\mathrm{g}},\;\hat{A}\right].
\end{equation}
The optimisation direction is specified through the Ax objective expressions, so the observations
themselves are not negated before being reported. This preserves the original physical meaning of
the objective values in the returned Pareto set.

\subsection{Output of Multi-Suction Cup Gripper Optimization}
\label{sec:multi_sc_output}

In addition to the common output described in Section~\ref{sec:common_output}, the multi-suction cup
optimization returns:
\begin{itemize}
    \item The Pareto-optimal set $\mathcal{P}^*$ in the reduced decision space, with corresponding
          objective values
    \item The associated radial gripper parameterizations
          $\mathbf{q} = [r,\alpha_1,\alpha_2,\alpha_3]$
\end{itemize}

These solutions represent optimal trade-offs among grasp diversity, the number of valid grasps, and
triangle area, corresponding to grasp versatility, robustness, and stability.

% ============================================================================
\section{Implementation of Two-Finger Gripper Optimization}
\label{sec:two_finger_gripper_optimization}

This section describes the problem-specific implementation for a two-finger gripper with a
spline-based gripping contour. The common framework components are described in
Section~\ref{sec:gripper_optimization_framework}.

The evaluation service constructs the gripper geometry, voxelizes the gripper and workpieces, and
calculates the grasp-quality objectives using a modified version of
Snuggle-Pack~\cite{snuggle_pack}. The same decision variables, parameter bounds,
geometry-generation pipeline, workpiece, orientations, and grasp-quality objectives are used for
both optimization approaches. This enables their performance to be compared under a common
evaluation budget and evaluation procedure.

\subsection{Decision Variables and Bounds for Two-Finger Gripper Problem}
\label{sec:two_finger_decision_variables}

A candidate gripper is represented by the decision vector introduced in
Section~\ref{subsec:methodology_decision_variables},
\begin{equation}
    \mathbf{q}
    =
    \begin{bmatrix}
        \ell_{\mathrm{gn}} &
        h_1 &
        h_2 &
        h_3
    \end{bmatrix}^{\mathrm{T}},
    \qquad d = 4,
\end{equation}
where the variables are transmitted to the evaluation API in the order
\begin{equation}
    \mathbf{q}
    =
    [
        \texttt{length\_gn},\;
        \texttt{h\_p1\_spline},\;
        \texttt{h\_p2\_spline},\;
        \texttt{h\_p3\_spline}
    ]^{\mathrm{T}}.
\end{equation}

The decision variables and their bounds are summarized in Table~\ref{tab:two_finger_variables}.

\begin{table}[htbp]
\centering
\caption{Decision variables of the two-finger gripper optimization problem.}
\label{tab:two_finger_variables}
\begin{tabularx}{\textwidth}{@{}l l l >{\raggedright\arraybackslash}X@{}}
\toprule
\textbf{Variable} & \textbf{Configuration parameter} & \textbf{Lower bound} & \textbf{Upper bound}
\\
\midrule
$\ell_{\mathrm{gn}}$ & \texttt{length\_gn}      & 20 & 50 \\
$h_1$                & \texttt{h\_p1\_spline}   & 2  & 10 \\
$h_2$                & \texttt{h\_p2\_spline}   & 2  & 10 \\
$h_3$                & \texttt{h\_p3\_spline}   & 2  & 10 \\
\bottomrule
\end{tabularx}
\end{table}

\nomenclature{$\ell_{\mathrm{gn}}$}{Length of the spline-based gripping contour}
\nomenclature{$h_j$}{Spline depth parameter at relative contour position $j/4$, $j=1,2,3$}

All four decision variables are treated as continuous. Candidate values generated by either
optimization framework are restricted to the bounded search space
\begin{equation}
    \mathcal{Q}
    =
    [20,50]
    \times
    [2,10]^3 \subset \mathbb{R}^4 .
\end{equation}

For evolutionary optimization, the variables are initially sampled within these bounds and
subsequently modified by the operators of the selected algorithm. For Bayesian optimization, Ax
initially evaluates the center of the search space and generates additional initialization
candidates using Sobol sampling. After sufficient observations have been collected, Ax transitions
to model-based candidate generation.

\subsection{Objective Function Evaluation of Two-Finger Gripper via API}
\label{sec:two_finger_api_evaluation}

For every candidate, the decision variables are sent to the evaluation API as an ordered list. The
API inserts these values into a complete gripper configuration, from which the jaw geometry is
generated:
\begin{equation}
    \mathbf{q}
    \;\longrightarrow\;
    \mathbf{c}_{\mathrm{gr}}(\mathbf{q})
    \;\longrightarrow\;
    \mathcal{G}(\mathbf{q}),
\end{equation}
where $\mathbf{c}_{\mathrm{gr}}(\cdot)$ denotes the mapping to the complete gripper configuration
and $\mathcal{G}(\cdot)$ the resulting jaw geometry.
\nomenclature{$\mathbf{c}_{\mathrm{gr}}(\cdot)$}{Complete gripper configuration generated from the decision vector}
\nomenclature{$\mathcal{G}(\cdot)$}{Gripper jaw geometry generated from the decision vector}

Besides the decision variables, the configuration $\mathbf{c}_{\mathrm{gr}}(\cdot)$ contains fixed parameters, including the contour
type, the basic jaw dimensions (depth, width, and length), the jaw angle, the stroke, the axial
offset of the contour, and the tool-center-point translation.

For every candidate, the evaluation request performs the following steps:
\begin{enumerate}
    \item Send the candidate decision vector to the external evaluation API.
    \item Map the decision variables to the complete gripper configuration
          $\mathbf{c}_{\mathrm{gr}}(\mathbf{q})$.
    \item Generate the spline-based gripper geometry using CadQuery.
    \item Voxelize the gripper and the configured workpiece in all predefined orientations.
    \item Calculate the FFT-based collision maps using the modified Snuggle-Pack procedure.
    \item Calculate one proximity KPI for each predefined orientation of the configured workpiece.
    \item Return the KPI vector to the selected optimization framework.
\end{enumerate}

\subsubsection{Parametrization of the Jaw Contour}
\label{sec:jaw_contour_parametrization}

The contact contour of the rigid jaw is described by an interpolating spline in the contour plane,
which is spanned by the jaw depth direction and the jaw length direction. The decision vector
consists of the contour length $\ell_{\mathrm{gn}}$ and three depth parameters $h_1$, $h_2$, and
$h_3$. These four values define the five control points of the spline,
\begin{equation}
    \mathbf{s}_0 = (0,\,0),
    \qquad
    \mathbf{s}_j
    =
    \left(
        h_j,\;
        \tfrac{j}{4}\,\ell_{\mathrm{gn}}
    \right),
    \quad j = 1,2,3,
    \qquad
    \mathbf{s}_4 = (0,\,\ell_{\mathrm{gn}}),
\end{equation}
so that the control points are distributed at the relative contour positions
$0,\ 0.25,\ 0.50,\ 0.75$, and $1.0$ along the contour length. The first and last control point are
pinned to zero depth, which guarantees that the contour begins and ends flush with the planar jaw
face and that the resulting profile is a closed recess rather than an open step. The intermediate
depths $h_1$, $h_2$, and $h_3$ control how far the contour intrudes into the jaw at the quarter,
half, and three-quarter positions and therefore determine the shape of the form fit. The contour
curve is the interpolating spline
\begin{equation}
    \begin{gathered}
        \mathcal{S}(\mathbf{q})
        =
        \operatorname{spline}
        \left(
            \mathbf{s}_0,\ldots,\mathbf{s}_4
        \right), \\[4pt]
        \mathbf{s}_j \in \mathbb{R}^{2},\; j = 0,\ldots,4,\quad
        \mathcal{S}(\mathbf{q}) : [0,1] \rightarrow \mathbb{R}^{2},\quad
        \mathcal{S}(\mathbf{q}) \in C^{1}\!\left([0,1],\mathbb{R}^{2}\right),\\[2pt]
        \exists\; 0 = u_0 < u_1 < \cdots < u_4 = 1 :\;
        \mathcal{S}(\mathbf{q})(u_j) = \mathbf{s}_j
        \;\text{(parameter values not necessarily } j/4)
    \end{gathered}
\end{equation}
which passes through all five control points.
\nomenclature{$\mathbf{s}_j$}{Spline control point $j$ of the jaw contour}
\nomenclature{$\mathcal{S}(\cdot)$}{Interpolating spline describing the jaw contact contour}

The spline is closed into a planar profile by straight segments up to a fixed cutting depth and is
then extruded through the full jaw depth to remove the enclosed material from the basic jaw body.
The geometry is constructed with CadQuery~\cite{cadquery2026} and exported to the mesh
representation required for evaluation. Since only $\ell_{\mathrm{gn}}$ and the three depth
parameters are varied, every candidate geometry differs from the others exclusively in the contact
contour, while all remaining jaw dimensions remain identical.

Each candidate geometry $\mathcal{G}(\mathbf{q})$ is generated and voxelized only once per
evaluation request and is subsequently reused for all configured grasps, so that the cost of the
geometry generation is independent of the number of objectives.

\subsubsection{Grasp Set Definition}
\label{sec:grasp_set}

A grasp specifies the pose in which the workpiece is presented to the jaws. In the evaluation
pipeline, the gripper remains stationary in its nominal pose, and the workpiece is reoriented for
each grasp.

Only the orientation component of the generated grasp poses is used, since the relative
gripper--workpiece position is determined by the collision- and proximity-based placement search
described in Section~\ref{sec:grasp_eval}. A grasp $\gamma_i$ is therefore fully described by a
spatial rotation $\mathbf{R}_i \in \mathrm{SO}(3)$ applied to the workpiece,
\begin{equation}
    \mathcal{W}_i = \mathbf{R}_i\,\mathcal{W},
    \qquad
    i = 1,\ldots,n_{\mathrm{or}},
\end{equation}
where $\mathcal{W}$ denotes the workpiece in its reference orientation. The translation components
of the generated poses are discarded. The grasp set is
\begin{equation}
    \mathcal{R}
    =
    \left\{
        \mathbf{R}_1,\ldots,\mathbf{R}_{n_{\mathrm{or}}}
    \right\},
    \qquad
    n_{\mathrm{or}} = 10 .
\end{equation}
\nomenclature{$\gamma_i$}{Grasp $i$ of the predefined grasp set}
\nomenclature{$\mathcal{R}$}{Set of predefined workpiece orientations of a two-finger optimization run}
\nomenclature{$n_{\mathrm{or}}$}{Number of predefined grasp orientations of a two-finger optimization run}
\nomenclature{$\mathcal{W}_i$}{Workpiece geometry in the orientation of grasp $\gamma_i$}

In the present study, $n_{\mathrm{or}} = 10$ orientations are configured.
Because the retained poses originate from the feasibility filtering and the farthest-point
clustering of the generation procedure, the orientations are neither restricted to a single rotation
axis nor obtained as independent samples of a parametric distribution. They are instead a spatially
spread subset of the physically admissible grasps of the given workpiece, which ensures that the
evaluated orientations are mutually dissimilar and at the same time realizable on the part geometry.

The grasp set is generated once before the optimization starts and is stored as part of the workpiece list. It is therefore identical for all
candidates and for all repetitions of a run, which makes the objective values of different candidates directly comparable. The orientations are fixed
quantities and are not optimization variables.

Each optimization run is defined for exactly one workpiece together with its corresponding
predefined grasp set. The workpiece geometry and the grasp set $\mathcal{R}$ are constant throughout
a run, and all objectives of a run refer to the same workpiece evaluated in its $n_{\mathrm{or}}$
predefined orientations. Different workpieces are treated as separate optimization problems and are
optimized in separate runs, since each workpiece requires its own task-specific jaw contour. 

\subsubsection{Grasp Evaluation}
\label{sec:grasp_eval}

The generated gripper is evaluated separately for every grasp $\gamma_i$,
$i = 1,\ldots,n_{\mathrm{or}}$, of the workpiece configured for the run. The gripper remains in
its assembly pose, while the workpiece is rotated into the orientation of the respective grasp.
Both geometries are voxelized into three-dimensional occupancy grids. For the gripper, the surface
voxelization of each part is closed by morphological dilation and hole filling. Every voxel
carries the label of the part it belongs to, so the two jaws can be addressed individually.

The evaluation builds on a modified version of Snuggle-Pack~\cite{snuggle_pack}, a
three-dimensional packing approach that combines FFT-based spatial analysis with proximity
heuristics. Instead of placing objects inside a container, the same principle is used to search
for the relative placement at which both jaws enclose the workpiece as closely as possible without
penetrating it.

For each jaw $j \in \{\mathrm{L},\mathrm{R}\}$ with occupancy grid $\chi_{\mathcal{J}_j}$, the
overlap with the workpiece grid $\chi_{\mathcal{W}_i}$ is computed for all discrete displacements
$\mathbf{v}$ in a single transform step:
\begin{equation}
    \zeta_{i,j}
    =
    \operatorname{DFT}^{-1}\!\left(
        \operatorname{DFT}(\chi_{\mathcal{J}_j})\,
        \operatorname{DFT}(\chi_{\mathcal{W}_i})
    \right),
    \qquad
    \zeta_{i,j}(\mathbf{v}) < \varepsilon_{\mathrm{col}}
    \;\Longleftrightarrow\;
    \mathbf{v} \text{ is collision-free}.
\end{equation}
The threshold $\varepsilon_{\mathrm{col}}$ absorbs the round-off of the single-precision FFT.
Because the grids are padded only to a common size, the transform is cyclic. Only the jaws enter
the collision test while the gripper base is not considered.

The overlap map of each jaw is then convolved with a one-sided Gaussian proximity kernel that
points along the closing direction of that jaw. The resulting proximity map is high for
displacements at which workpiece material lies directly in front of the jaw face, that is, where
a small further closing motion would establish contact. The kernel's standard deviations are given
in millimetres and converted to voxel units via the voxel edge length. The kernel acts only in the
plane orthogonal to the gripper length axis, so in practice only the width and height deviations
shape the score. Colliding displacements are excluded from the proximity map.

The placement search follows the kinematics of a parallel gripper. For every position in the
length--height plane, each jaw independently selects its best collision-free position along the
closing axis. The two jaw proximities are then summed and maximized over that plane, which yields
the proximity score $p_i$. Both jaws therefore share a common placement in length and height but
move independently along the closing direction. All operations run as GPU-accelerated PyTorch FFT
convolutions. The score is dimensionless, unnormalized and Gaussian-weighted. It increases the more
closely both jaws enclose the workpiece, but it is neither a distance nor a voxel count, and its
magnitude depends on the workpiece and on the voxel resolution.

The difference between the two selected jaw positions along the closing axis defines a signed
opening width $\delta_i$. A negative value means the jaws close onto the workpiece. A non-negative
value means the placement would require the jaws to open beyond their modelled position. Only the
sign is evaluated; the magnitude is not compared against the stroke. The KPI reported to the
optimization framework is
\begin{equation}
    k_i(\mathbf{q}) =
    \begin{cases}
        p_i(\mathbf{q}), & \delta_i(\mathbf{q}) < 0,\\[2pt]
        0,               & \delta_i(\mathbf{q}) \geq 0,
    \end{cases}
    \qquad i = 1,\ldots,n_{\mathrm{or}}.
\end{equation}
Feasibility is assessed per grasp, so an infeasible orientation penalizes only its own objective.
Since collision-free proximity values are non-negative, zero is the worst attainable outcome. An
infeasible orientation is therefore treated as a worst-case objective value rather than as a hard
constraint.

\nomenclature{$\chi_{\mathcal{J}_j}$}{Voxel occupancy grid of gripper jaw $j$}
\nomenclature{$\chi_{\mathcal{W}_i}$}{Voxel occupancy grid of the workpiece in grasp orientation $i$}
\nomenclature{$\zeta_{i,j}$}{Overlap map of jaw $j$ and the workpiece for grasp $\gamma_i$}
\nomenclature{$\mathbf{v}$}{Relative displacement between the voxel grids}
\nomenclature{$\varepsilon_{\mathrm{col}}$}{Numerical collision threshold of the FFT-based overlap map}
\nomenclature{$\operatorname{DFT}$}{Discrete Fourier transform}
\nomenclature{$p_i$}{Maximized jaw proximity score of grasp $\gamma_i$}
\nomenclature{$\delta_i$}{Signed opening width of the placement selected for grasp $\gamma_i$}
\nomenclature{$k_i$}{Proximity KPI \texttt{max\_prox} of grasp $\gamma_i$ reported to the optimizer}

\subsection{Objective Evaluation of Two-Finger Gripper}
\label{sec:two_finger_objective_evaluation}

The evaluation service returns one grasp-quality KPI per configured grasp. Because a single
workpiece and its associated grasp set are used per optimization run, the number of objectives is
determined directly by the size of that grasp set:
\begin{equation}
    m = n_{\mathrm{or}} .
\end{equation}
The individual KPIs are not aggregated into a scalar value; each grasp retains its own objective,
which preserves the information about which orientations a given contour favours.

Since a larger proximity KPI represents a better grasp result, the natural optimization problem is
\begin{equation}
    \begin{gathered}
        \max_{\mathbf{q}\in\mathcal{Q}} k_i(\mathbf{q}),
        \qquad
        i = 1,\ldots,m , \\[4pt]
        \mathcal{Q} \subset \mathbb{R}^{4}
        \;,\quad
        m = n_{\mathrm{or}} = 10,\\[2pt]
        k_i : \mathcal{Q} \rightarrow \mathcal{K} \subseteq \mathbb{R}_{\geq 0}
    \end{gathered}
\end{equation}
with the objective vector
\begin{equation}
    \mathbf{f}(\mathbf{q})
    =
    \begin{bmatrix}
        k_1(\mathbf{q}) &
        k_2(\mathbf{q}) &
        \cdots &
        k_{m}(\mathbf{q})
    \end{bmatrix}^{\mathrm{T}}
    \in \mathbb{R}_{\geq 0}^{m}.
\end{equation}
\nomenclature{$k_i(\cdot)$}{Grasp-quality KPI (\texttt{max\_prox}) of grasp $\gamma_i$}

Because the four decision variables define a single jaw contour that is shared by all grasps, the
objectives are generally conflicting: a contour whose depth profile conforms closely to the
workpiece in one orientation does not necessarily conform in the remaining ones. The optimization
therefore yields Pareto-optimal contours that represent different compromises in geometric
conformity across the ten predefined orientations.

Following the convention of Section~\ref{sec:objective_direction_conventions}, the evolutionary
algorithms receive the negated KPI vector
\begin{equation}
    \mathbf{f}_{\mathrm{EA}}(\mathbf{q})
    =
    -\,\mathbf{f}(\mathbf{q})
    =
    \begin{bmatrix}
        -k_1(\mathbf{q}) &
        -k_2(\mathbf{q}) &
        \cdots &
        -k_{m}(\mathbf{q})
    \end{bmatrix}^{\mathrm{T}},
\end{equation}
whereas the Bayesian implementation uses the unchanged observations
\begin{equation}
    \mathbf{f}_{\mathrm{BO}}(\mathbf{q})
    =
    \mathbf{f}(\mathbf{q}).
\end{equation}
Consequently, the observed KPI values do not need to be negated before they are attached to the
corresponding Ax trial, which preserves the physical meaning of the KPI values when the final Pareto
frontier is returned. For performance comparisons, the results of both optimization frameworks are
converted to the common natural maximization convention $\mathbf{f}(\mathbf{q})$, so that all
reported objective values correspond to the ten grasp-quality KPIs in their original direction and
magnitude.

\subsection{Output of Two-Finger Gripper Optimization}
\label{sec:two_finger_output}

In addition to the common output described in Section~\ref{sec:common_output}, the two-finger
optimization returns:
\begin{itemize}
    \item The Pareto-optimal decision vectors
          $\mathbf{q}=[\ell_{\mathrm{gn}},h_1,h_2,h_3]^{\mathrm{T}}$
    \item The corresponding objective KPI values $\mathbf{f}(\mathbf{q})$
\end{itemize}

The non-dominated solutions therefore represent the best trade-offs between grasp quality across
the predefined orientations of the configured workpiece.